% from ega0-19-fr.tex
\oldpage[0\textsubscript{IV-1}]{74}

\begin{proof}
En vertu de (19.1.7), on a des diagrammes commutatifs
\[
  \xymatrix{
    0\ar[r] & M_n\ar[r]^{u_n} & N_n\ar[r]^{p_n} & P_n\ar[r] & 0\\
    0\ar[r] & M_{n+1}\ar[u]^f\ar[r]_{u_{n+1}} &
      N_{n+1}\ar[u]^g\ar[r]_{p_{n+1}} & P_{n+1}\ar[u]^h\ar[r] & 0
  }
\]
où les lignes sont des suites exactes \emph{scindées}, autrement dit, il existe
pour chaque $n$ un homomorphisme $s_n:P_n\to N_n$ tel que
$p_n\circ s_n=1_{P_n}$. Nous allons montrer que l'on peut, par récurrence sur
$n$, définir un homomorphisme $s'_n:P_n\to N_n$ tel que
$p_n\circ s'_n=1_{P_n}$ et que les diagrammes
\[
  \xymatrix{
    P_n\ar[r]^{s'_n} & N_n\\
    P_{n+1}\ar[u]^h\ar[r]_{s'_{n+1}} & N_{n+1}\ar[u]^g
  }
\]
soient commutatifs. En effet, $g\circ s_{n+1}-s'_n\circ h$ est un
homomorphisme de $P_{n+1}$ dans
\[
  u_n(M_n)=u_{n+1}(M_{n+1})/
  \mathfrak J_nu_{n+1}(M_{n+1})~;
\]
comme $P_{n+1}$ est un $(A/\mathfrak J_{n+1})$-module \emph{projectif}, cet
homomorphisme se factorise en
\[
  P_{n+1}\xrightarrow{t_{n+1}}u_{n+1}(M_{n+1})
  \longrightarrow
  u_{n+1}(M_{n+1})/\mathfrak J_nu_{n+1}(M_{n+1}),
\]
d'où l'on conclut aussitôt que $s'_{n+1}=s_{n+1}-t_{n+1}$ répond à la
question. Cela étant, de la décomposition de $N_n$ en somme directe de
$u_n(M_n)$ et de $s'_n(P_n)$, on déduit aussitôt un homomorphisme
$w_n:N_n\to M_n$ inverse à gauche de $u_n$ et tel que les diagrammes
\[
  \xymatrix{
    N_n\ar[r]^{w_n} & M_n\\
    N_{n+1}\ar[u]^g\ar[r]_{w_{n+1}} & M_{n+1}\ar[u]^f
  }
\]
soient commutatifs. Le système projectif $(w_n)$ admet alors une limite
projective $\widehat w:\widehat N\to\widehat M=M$, d'où par composition avec
l'homomorphisme canonique $N\to\widehat N$, un homomorphisme $w:N\to M$ tel
que, pour tout $n$, l'endomorphisme $(w\circ u)_n$ de
$M_n=M/\mathfrak J_nM$ déduit de $w\circ u$ par passage aux quotients soit
l'identité~; cela entraîne que $w\circ u=1_M$ puisque $M$ est séparé et
complet.
\end{proof}

\begin{proposition}[19.1.9]
\label{0.19.1.9-fr}
Soient $A$ un anneau topologique préadmissible ($\mathbf{0}_{\mathrm I}$,
7.1.2), $\mathfrak Q$ un idéal de définition de $A$,
$(\mathfrak J_\lambda)$ un système fondamental d'idéaux ouverts de $A$.
Soient $M$, $N$ deux $A$-modules topologiques dont les topologies sont déduites
de celles de $A$, et tels que pour tout $\lambda$,
$N/\mathfrak J_\lambda N$ soit un $(A/\mathfrak J_\lambda)$-module projectif
(cf. (19.2.3)). Soit $u:M\to N$ un $A$-homomorphisme. Alors les conditions
suivantes sont équivalentes~:
\begin{enumerate}
  \item[a)] $u$ est formellement inversible à gauche.

  